\documentclass[11pt,letterpaper]{article}

\usepackage[utf8]{inputenc}
\usepackage[T1]{fontenc}
\usepackage[spanish,es-nodecimaldot]{babel}
\usepackage{lmodern}
\usepackage{geometry}
\usepackage{microtype}
\usepackage{booktabs}
\usepackage{tabularx}
\usepackage{array}
\usepackage{enumitem}
\usepackage{xcolor}
\usepackage{listings}
\usepackage{amsmath}
\usepackage{graphicx}
\usepackage{hyperref}

\geometry{margin=2.4cm}
\setlength{\parindent}{0pt}
\setlength{\parskip}{0.55em}
\setlist[itemize]{leftmargin=1.4em,itemsep=0.2em,topsep=0.3em}
\hypersetup{colorlinks=true,linkcolor=black,urlcolor=blue!55!black}

\definecolor{codegreen}{RGB}{0,105,70}
\definecolor{codegray}{RGB}{90,90,90}
\definecolor{codebackground}{RGB}{247,247,247}
\lstset{
  language=C,
  basicstyle=\ttfamily\small,
  backgroundcolor=\color{codebackground},
  keywordstyle=\color{blue!65!black},
  commentstyle=\color{codegreen},
  stringstyle=\color{red!60!black},
  numbers=left,
  numberstyle=\tiny\color{codegray},
  frame=single,
  breaklines=true,
  showstringspaces=false,
  tabsize=2
}

\newcolumntype{Y}{>{\raggedright\arraybackslash}X}

\title{\textbf{Proyecto 2 - Celeste}}
\author{Miguel Donis -- 22993\\Ian Farrington -- -----}
\date{\today}

\begin{document}

\maketitle
\thispagestyle{empty}
\vfill

\begin{center}
  \large Diseño e implementación de un videojuego embebido\\[0.5em]
  \normalsize STM32F446RE, ESP32, pantalla ILI9341 y audio por DAC
\end{center}

\vfill
\newpage

\tableofcontents
\newpage

\section{Resumen}

El proyecto implementa un videojuego de plataformas para dos jugadores inspirado en
\textit{Celeste Classic}. La lógica principal se ejecuta en una STM32F446RE y se
muestra en una pantalla ILI9341 de $240\times320$ píxeles. Dos ESP32 reciben los
controles por Bluetooth y transmiten el estado de cada jugador a la STM principal.
Una segunda STM32 genera la música mediante DAC, DMA y un temporizador.

Los niveles se almacenan como texto RGB565 en una tarjeta microSD. Al entrar a un
nivel, la STM compone el fondo y la geometría, guarda el resultado en una región
reservada de su Flash interna y crea un mapa compacto de colisiones en RAM. Esta
distribución evita leer la SD durante el movimiento y permite restaurar únicamente
la región ocupada por cada sprite.

\section{Arquitectura del sistema}

El sistema se dividió por función para que el procesamiento gráfico, los controles y
el audio no compitan por el mismo periférico ni bloqueen el ciclo del juego.

\begin{table}[h]
\centering
\caption{Funciones de los componentes principales.}
\begin{tabularx}{\textwidth}{@{}lY@{}}
\toprule
\textbf{Componente} & \textbf{Función} \\
\midrule
STM32 principal & Ejecuta el juego, carga niveles, calcula física y colisiones,
controla la ILI9341 y mantiene el puntaje. \\
ESP32 J1 y J2 & Reciben cada control Bluetooth con Bluepad32 y envían tramas UART
independientes. \\
STM32 de música & Recibe el estado del juego por SPI2 y genera las notas por DAC. \\
MicroSD & Almacena el fondo general, cinco niveles y tres imágenes finales. \\
ILI9341 & Presenta el escenario, los sprites, el marcador y los resultados. \\
Amplificador & Adapta la señal analógica del DAC para manejar la bocina. \\
\bottomrule
\end{tabularx}
\end{table}

La pantalla utiliza un bus paralelo de 8 bits. La SD se conecta a SPI1, los ESP32 a
UART5 y USART3, y la STM de música a SPI2. Esto evita compartir buses entre cargas
con requisitos distintos.

\section{Organización del programa principal}

El programa usa un ciclo principal sin sistema operativo. Cada tarea revisa el
tiempo con \texttt{HAL\_GetTick()} y se ejecuta cuando corresponde. La física se
actualiza cada 20 ms, equivalente a 50 actualizaciones por segundo. En cada vuelta se
realizan las siguientes operaciones:

\begin{enumerate}[leftmargin=1.6em,itemsep=0.2em]
  \item Actualizar la recepción de los dos controles.
  \item Revisar eventos de inicio y el estado de la pantalla actual.
  \item Actualizar la física de ambos jugadores cuando vence el periodo de 20 ms.
  \item Procesar peligro, salida, puntaje y cambio de nivel.
  \item Restaurar el fondo anterior y dibujar los sprites en su nueva posición.
  \item Reenviar periódicamente a la STM de música el tema que debe reproducir.
\end{enumerate}

El flujo general tiene cuatro estados visibles: menú, partida, resultados e imágenes
finales. \texttt{START} inicia la partida; después del último nivel muestra el
resultado y nuevas pulsaciones recorren las tres imágenes finales.

\section{Gráficos y administración de memoria}

\subsection{Formato de texturas}

Los colores se guardan en RGB565, por lo que cada píxel ocupa 16 bits. Los archivos
de la SD conservan el formato hexadecimal original y pueden contener dos píxeles
empaquetados en una palabra de 32 bits o valores RGB565 individuales. Una tabla de
niveles registra el nombre, altura jugable, posición inicial, color transparente y
formato de cada archivo. El número de niveles se calcula con \texttt{sizeof}, por lo
que la lógica de progresión no contiene un límite escrito manualmente.

La composición se hace una fila a la vez:

\begin{equation}
C(x,y)=
\begin{cases}
B(x,y), & L(x,y)=T,\\
L(x,y), & L(x,y)\neq T,
\end{cases}
\end{equation}

donde $B$ es el fondo general, $L$ es la textura del nivel y $T$ es su color
transparente. Debajo del área jugable se genera una franja de marcador de color
RGB565 \texttt{0x5E7B}.

\subsection{Caché del nivel activo}

Una imagen completa requiere

\begin{equation}
240\cdot320\cdot2=153\,600\ \text{bytes}.
\end{equation}

La STM32F446RE tiene 128 KiB de RAM, por lo que el cuadro completo no cabe junto con
el resto del programa. Se reservaron los sectores 6 y 7 de Flash, desde
\texttt{0x08040000}, como una caché de 256 KiB. El enlazador limita el firmware a los
primeros 256 KiB y evita que el código invada esa zona.

Durante la carga se leen el fondo y el nivel desde la SD, se compone cada fila, se
genera su colisión y se escribe el resultado en Flash. Al finalizar se compara un
hash FNV-1a de los datos producidos con el contenido almacenado. La SD queda fuera
del ciclo de juego hasta el siguiente cambio de nivel.

Para mover un jugador no se redibuja toda la pantalla. Primero se recupera desde la
Flash la región que ocupaba su sprite y después se dibuja el nuevo cuadro con
transparencia. Esto reduce transferencia a la ILI9341 y evita rastros visibles.

\section{Colisiones y lógica del escenario}

El mapa de colisiones utiliza dos bits por píxel. Para los 76,800 píxeles ocupa
19,200 bytes de RAM y admite cuatro tipos:

\begin{table}[h]
\centering
\caption{Tipos guardados en el mapa de colisiones.}
\begin{tabular}{@{}cll@{}}
\toprule
\textbf{Valor} & \textbf{Tipo} & \textbf{Comportamiento} \\
\midrule
0 & Libre & El jugador puede ocupar el píxel. \\
1 & Sólido & Suelo, pared, techo y franja del marcador. \\
2 & Peligro & Rojo \texttt{0xF800}; reinicia la posición del jugador. \\
3 & Salida & Amarillo \texttt{0xFF60} o \texttt{0xFFE0}; completa el nivel. \\
\bottomrule
\end{tabular}
\end{table}

Un píxel transparente es libre. Cualquier otro color de la textura es sólido, salvo
los colores reservados para peligro y salida. Los límites de la pantalla también se
consideran sólidos.

La caja de colisión del jugador mide $18\times14$ píxeles. Para verificar una
posición, el código recorre todo el rectángulo y consulta el mapa. El movimiento se
resuelve primero en el eje horizontal y luego en el vertical, avanzando un píxel por
iteración. Si el siguiente paso colisiona, la posición queda en el último punto
válido y la velocidad de ese eje se anula. Este recorrido evita atravesar superficies
cuando una velocidad representa varios píxeles por actualización.

El suelo no se almacena como un tipo diferente de la pared. Se reconoce probando la
misma caja un píxel debajo del jugador. Las paredes se buscan a un píxel de los
laterales y el techo se detecta cuando el movimiento vertical ascendente encuentra
una región sólida.

\section{Movimiento de los jugadores}

La posición y velocidad usan punto fijo con 256 unidades internas por píxel. Esto
permite aceleraciones suaves sin depender de números flotantes durante cada cuadro.
Los multiplicadores decimales se aplican al compilar para producir constantes
enteras redondeadas.

Cada jugador conserva su posición, velocidad, orientación, cuadro de animación,
contador de saltos y estados temporales. La velocidad horizontal se aproxima de
forma gradual a una velocidad objetivo. Se usan aceleraciones separadas para suelo y
aire, además de desaceleración cuando se libera la dirección.

El contador de saltos tiene tres estados: salto normal disponible, segundo salto
disponible y saltos agotados. El salto normal alcanza actualmente 2.5 veces la altura
del jugador y habilita el segundo. El segundo salto alcanza tres veces la altura y
depende de la dirección mantenida:

\begin{itemize}
  \item Sin dirección, o solamente arriba: impulso vertical.
  \item Dirección lateral sin arriba: desplazamiento horizontal con distancia
  limitada; la gravedad se pausa hasta alcanzar el límite o chocar.
  \item Arriba y una dirección lateral: impulso diagonal.
\end{itemize}

Al aterrizar se restablece el contador. El salto de pared expulsa al jugador hacia el
lado opuesto y aplica un bloqueo corto de dirección para conservar el impulso. El
contacto lateral durante una caída limita la velocidad al 55\% de la velocidad
terminal normal. Un salto de pared no recupera un segundo salto ya consumido.

También se mantiene un búfer de cuatro cuadros para el salto normal. Una pulsación
recibida ligeramente antes de cumplir la condición puede ejecutarse en las
actualizaciones siguientes, lo que hace el control menos dependiente del instante
exacto de recepción.

\section{Controles inalámbricos}

Cada ESP32 ejecuta Bluepad32 y acepta únicamente la dirección Bluetooth registrada
para su jugador. El estado se convierte en una máscara de siete bits: izquierda,
derecha, arriba, abajo, Cruz, Círculo y Start. Las combinaciones se conservan porque
cada bit se procesa de manera independiente.

El ESP32 transmite por UART a 115200 bit/s una trama cada 20 ms:

\begin{center}
\texttt{[0xA5, ID, botones, estado, CRC-8]}
\end{center}

El byte de estado es \texttt{0xA1} cuando el control está conectado y
\texttt{0xA0} cuando está desconectado. El CRC-8 usa polinomio \texttt{0x07} y valor
inicial \texttt{0x00}. La STM reconstruye las tramas mediante interrupciones de un
byte, valida sincronización, ID, estado y CRC, y mantiene una cola independiente para
cada jugador.

Se distinguen botones mantenidos de eventos de pulsación. Las direcciones se leen
como estado continuo; salto, segundo salto y Start generan un evento únicamente al
pasar de liberado a presionado. Si no llega una trama válida durante 300 ms, las
acciones mantenidas se liberan para evitar que el jugador continúe moviéndose por una
desconexión.

\begin{table}[h]
\centering
\caption{Asignación de controles.}
\begin{tabular}{@{}ll@{}}
\toprule
\textbf{Entrada} & \textbf{Acción} \\
\midrule
Cruceta izquierda/derecha & Movimiento horizontal \\
Cruceta arriba & Dirección del segundo salto \\
Cruceta abajo & Disponible en el protocolo \\
Cruz & Salto normal o salto de pared \\
Círculo & Segundo salto direccional \\
Start & Iniciar partida o avanzar pantallas finales \\
\bottomrule
\end{tabular}
\end{table}

\section{Sistema de sonido}

La STM principal envía un byte por SPI2: \texttt{'M'} para menú y resultados, o
\texttt{'T'} para la partida. El comando se manda al cambiar de estado y se repite
cada segundo para recuperar la sincronización si una de las placas se reinicia. La
STM de música funciona como esclavo SPI en modo 0, con datos de 8 bits y bit más
significativo primero.

Las canciones se almacenan como dos arreglos paralelos: frecuencia y duración de
cada nota. Una máquina de estados basada en \texttt{HAL\_GetTick()} avanza la canción
sin usar \texttt{HAL\_Delay()}, por lo que la recepción SPI permanece activa. Un
comando repetido no reinicia la melodía.

El audio se genera con una tabla senoidal de 128 muestras. TIM6 produce el disparo
del DAC y su periodo se calcula para cada frecuencia $f$:

\begin{equation}
ARR \approx \frac{84\,000\,000}{128f}-1.
\end{equation}

DMA transfiere la tabla al DAC1 en modo circular. Una frecuencia de cero detiene el
temporizador y coloca el DAC en 2048, el punto medio de 12 bits. El volumen digital
se ajusta al 50\% escalando la amplitud respecto a ese punto medio.

\subsection{Circuito de salida}

La salida \texttt{PA4/DAC1\_OUT1} se conecta a un capacitor electrolítico de
$10\,\mu\text{F}$, con el terminal positivo hacia la STM. El otro terminal entra a
un canal del amplificador. El capacitor elimina el nivel DC de la señal. La bocina se
conecta entre las dos terminales de salida del mismo canal del amplificador; ninguna
terminal de la bocina se conecta a tierra.

\begin{center}
\texttt{PA4 DAC $\rightarrow$ 10 uF $\rightarrow$ entrada L del amplificador
$\rightarrow$ L+/L- $\rightarrow$ bocina}
\end{center}

El amplificador trabaja con 5 V y comparte tierra con ambas STM y los ESP32. Se usa
una fuente común regulada para evitar que los picos de consumo del amplificador o de
los radios Bluetooth afecten la lógica digital.

\section{Puntaje y final de partida}

Cada nivel muestra una franja inferior con el identificador y puntaje de ambos
jugadores. Cuando alguien toca la salida amarilla, el primer evento procesado recibe
un punto, ambos jugadores se reinician y se carga el nivel siguiente. Los puntajes se
mantienen en variables de 16 bits durante toda la partida.

Después del quinto nivel se compara el total y se muestra ganador o empate sobre la
pantalla de inicio. La música cambia al tema de menú. A continuación, cada pulsación
de Start carga desde la SD una de las tres imágenes finales. Al terminar la secuencia
se regresa al menú y una partida nueva reinicia ambos puntajes.

\section{Esquemáticos}

Esta sección resume las conexiones físicas empleadas en el proyecto. Las señales
lógicas trabajan a 3.3 V y todos los módulos comparten la misma referencia de
tierra. Los diagramas esquemáticos se agregarán junto a estas tablas.

\subsection{Pantalla ILI9341}

La pantalla utiliza una interfaz paralela de 8 bits. Las líneas de datos y control
se conectan a la STM principal de la siguiente manera:

\begin{table}[htbp]
\centering
\small
\caption{Conexiones entre la STM principal y la pantalla ILI9341.}
\begin{tabularx}{0.78\textwidth}{@{}YY@{}}
\toprule
\textbf{Señal ILI9341} & \textbf{Pin STM principal} \\
\midrule
D0 & PA9 \\
D1 & PC7 \\
D2 & PA10 \\
D3 & PB3 \\
D4 & PB5 \\
D5 & PB4 \\
D6 & PB10 \\
D7 & PA8 \\
RD & PA0 \\
WR & PA1 \\
RS/DC & PA4 \\
RST & PC1 \\
CS & PB0 \\
VCC y GND & 3.3 V y tierra común \\
\bottomrule
\end{tabularx}
\end{table}

\subsection{Tarjeta microSD}

La tarjeta se comunica mediante SPI1 y se utiliza únicamente durante la carga de
un nivel o de una imagen final.

\begin{table}[htbp]
\centering
\small
\caption{Conexiones del módulo microSD.}
\begin{tabularx}{0.78\textwidth}{@{}YY@{}}
\toprule
\textbf{Señal microSD} & \textbf{Pin STM principal} \\
\midrule
SCK & PA5 / SPI1\_SCK \\
MISO & PA6 / SPI1\_MISO \\
MOSI & PA7 / SPI1\_MOSI \\
CS/SS & PB6 \\
VCC y GND & 3.3 V y tierra común \\
\bottomrule
\end{tabularx}
\end{table}

\subsection{Controles inalámbricos}

Cada ESP32 envía por UART únicamente la trama correspondiente a su jugador. No se
utiliza la línea de recepción de los ESP32.

\begin{table}[htbp]
\centering
\small
\caption{Conexiones de los ESP32 con la STM principal.}
\begin{tabularx}{\textwidth}{@{}l l Y@{}}
\toprule
\textbf{Módulo} & \textbf{STM principal} & \textbf{Función} \\
\midrule
ESP32 J1 GPIO17/TX2 & PD2/UART5\_RX & Recepción del control del jugador 1 \\
ESP32 J2 GPIO17/TX2 & PC11/USART3\_RX & Recepción del control del jugador 2 \\
GND de ambos ESP32 & GND común & Referencia de las señales UART \\
\bottomrule
\end{tabularx}
\end{table}

\subsection{STM de música y salida de audio}

La STM principal actúa como maestro de SPI2 y envía un byte para seleccionar el
tema. La STM de música funciona como esclavo y solamente recibe, por lo que no se
requiere una línea MISO.

\begin{table}[htbp]
\centering
\small
\caption{Comunicación entre la STM principal y la STM de música.}
\begin{tabularx}{\textwidth}{@{}l l Y@{}}
\toprule
\textbf{STM principal} & \textbf{STM de música} & \textbf{Señal} \\
\midrule
PC3/SPI2\_MOSI & PC1/SPI2\_MOSI & Comando de selección de tema \\
PB13/SPI2\_SCK & PB13/SPI2\_SCK & Reloj SPI \\
PB9/SPI2\_NSS & PB12/SPI2\_NSS & Selección del esclavo \\
GND & GND & Tierra común \\
\bottomrule
\end{tabularx}
\end{table}

La salida \texttt{PA4/DAC1\_OUT1} de la STM de música se conecta al terminal
positivo de un capacitor electrolítico de $10\,\mu\text{F}$. El terminal negativo
entra al canal izquierdo del amplificador. La bocina se conecta entre
\texttt{L+} y \texttt{L-}; ninguna de sus terminales se conecta directamente a
tierra.

\subsection{Alimentación}

La alimentación se distribuye en paralelo desde una fuente regulada de 5 V. Las
dos placas Nucleo reciben esta tensión por \texttt{E5V}, los ESP32 por
\texttt{VIN/5V} y el amplificador por su entrada de 5 V. La pantalla y la microSD
utilizan 3.3 V. Todas las tierras se conectan entre sí.

\begin{figure}[htbp]
\centering
% Sustituir el recuadro por, por ejemplo:
% \includegraphics[width=\textwidth]{figuras/esquematico_general.pdf}
\fbox{%
  \parbox[c][5cm][c]{0.88\textwidth}{%
    \centering
    Espacio para colocar los esquemáticos del proyecto
  }%
}
\caption{Esquemático general de conexiones del sistema.}
\label{fig:esquematico-general}
\end{figure}

\section{Análisis de la implementación}

La principal restricción fue la memoria. Leer directamente desde la SD durante el
juego producía pausas y hacía que la restauración del personaje dependiera del tiempo
de acceso del sistema de archivos. La caché en Flash resolvió este problema sin
reservar un cuadro imposible de alojar en RAM. Como costo, cada cambio de nivel
requiere borrar sectores internos y esperar la carga completa. El desgaste se limita
porque la Flash no se reescribe al morir, únicamente al cambiar de nivel.

El mapa de dos bits por píxel permite consultar colisiones en tiempo constante y
distinguir geometría, peligro y salida con un uso moderado de RAM. Separar posición
subpíxel y colisión por píxel mantiene movimiento suave sin complicar la geometría.

La división del sonido en otra STM evita que la temporización del DAC y el acceso DMA
afecten la pantalla o la física. Del mismo modo, cada ESP32 tiene su propia UART, por
lo que los dos jugadores no comparten una línea ni requieren arbitraje.

Entre las limitaciones actuales, una llegada simultánea a la salida favorece al
primer jugador procesado, las texturas en texto requieren más tiempo de carga que un
formato binario y la comunicación de música no incluye confirmación de recepción.
Estas decisiones son aceptables para la escala actual y mantienen el código legible.

\section{Uso de herramientas de inteligencia artificial}

Durante el desarrollo se utilizó \textbf{ChatGPT-5.6 Sol} como herramienta de apoyo
para analizar problemas relacionados con la física, las colisiones y la organización
del código. También se utilizó para revisar la redacción de este informe. Las
decisiones de diseño, la integración del hardware, la adaptación del código y las
pruebas fueron realizadas y verificadas por los autores.

\subsection{Prompts utilizados}

Algunas de las consultas realizadas durante el proyecto fueron las siguientes:

\begin{enumerate}[leftmargin=1.6em,itemsep=0.35em]
  \item ``Ayúdame a ajustar la gravedad, la altura de los saltos y el movimiento
  horizontal para que el personaje se mueva de forma fluida.''
  \item ``Revisa la lógica de colisiones por píxel y explica cómo evitar que el
  jugador atraviese paredes, suelo o techo.''
  \item ``Revisa el informe realizado y adáptalo evitando faltas de ortografía y
  redundancias para tener un resultado más profesional.''
  \item ``En base a los comentarios realizados hasta el momento, completa los
  comentarios del código manteniendo buenas prácticas para que sea fácil de
  entender.''
\end{enumerate}

\subsection{Ajuste de la física}

El siguiente fragmento resume el uso de punto fijo y multiplicadores para ajustar el
movimiento sin ejecutar operaciones de punto flotante en cada cuadro. Este bloque se
desarrolló con apoyo de IA y después se calibró mediante pruebas en el juego.

\begin{lstlisting}[caption={Parámetros principales de movimiento ajustados con apoyo de IA.}]
#define GAME_FIXED_POINT_ONE 256L
#define GAME_JUMP_HEIGHT_MULTIPLIER 2.5f
#define GAME_SECOND_JUMP_HEIGHT_MULTIPLIER 3.0f
#define GAME_GRAVITY_MULTIPLIER 2.0f

#define GAME_SCALE_ROUNDED(base, multiplier) \
  ((int32_t)(((float)(base) * (multiplier)) + 0.5f))

gravity = (velocity_y >= -GAME_APEX_SPEED &&
           velocity_y <= GAME_APEX_SPEED)
              ? GAME_APEX_GRAVITY
              : GAME_GRAVITY;
velocity_y = Game_Approach(velocity_y,
                           GAME_MAX_FALL_SPEED,
                           gravity);
\end{lstlisting}

La solución separa la gravedad normal de la gravedad cercana al punto más alto del
salto. Los multiplicadores permiten modificar alturas y aceleración con valores
decimales, mientras que el resultado final se convierte a las unidades enteras del
sistema de punto fijo. Esto permitió ajustar el movimiento sin introducir cambios
bruscos entre cuadros.

\subsection{Colisiones por píxel}

Para impedir que una velocidad alta atraviese una superficie, se revisa el
rectángulo completo del jugador y el desplazamiento se aplica un píxel a la vez. El
siguiente bloque también fue desarrollado con apoyo de IA y adaptado al mapa de
colisiones del proyecto.

\begin{lstlisting}[caption={Comprobación de la caja completa del jugador.}]
static uint8_t Game_CollidesAt(int32_t x, int32_t y)
{
  for (uint16_t row = 0; row < GAME_PLAYER_HEIGHT; row++)
  {
    for (uint16_t column = 0; column < GAME_PLAYER_WIDTH; column++)
    {
      if (Game_IsSolidPixel(x + column, y + row))
      {
        return 1U;
      }
    }
  }
  return 0U;
}

while (current_x != target_x)
{
  int32_t next_x = current_x + ((target_x > current_x) ? 1L : -1L);
  if (Game_CollidesAt(next_x, player_y))
  {
    velocity_x = 0L;
    break;
  }
  current_x = next_x;
}
\end{lstlisting}

La caja completa evita detectar únicamente una esquina del sprite. El recorrido
intermedio garantiza que cada posición sea validada antes de avanzar. El mismo
criterio se aplica en el eje vertical, por lo que suelo, techo y paredes se resuelven
con una sola función de colisión.

\section{Conclusiones}

El proyecto integra gráficos, almacenamiento, comunicación inalámbrica, física y
audio en una plataforma con recursos limitados. La separación entre imagen compuesta
en Flash y mapa de colisiones en RAM permitió mantener una actualización estable sin
consultar la SD durante la partida. La lógica por estados y las tareas basadas en
tiempo evitan bloqueos entre controles, animación y música.

La implementación queda preparada para agregar niveles mediante su tabla de
configuración y para ajustar el movimiento mediante multiplicadores definidos en el
código. La arquitectura de dos controles y dos procesadores STM también mantiene
separadas las funciones de juego y reproducción de audio.

\end{document}
